\documentclass[12pt]{article}
\usepackage[T1]{fontenc}
\usepackage[utf8]{inputenc}
\usepackage{lmodern}
\usepackage{textcomp}
\usepackage{enumitem}
\usepackage{geometry}
\usepackage{graphicx}
\usepackage{amsmath, amssymb}
\geometry{margin=1in}
\begin{document}
\begin{center}
Mathematics - Undergraduate Level Assessment 25 \
Total Marks: 40 \
Answer all questions.
\end{center}

\section*{Multiple Choice (10 marks)}
\begin{enumerate}[label=\arabic*.]
\item Let x\_1 = 1 and x\_(n+1) = sqrt(3+2 x\_n) for all positive integers n. If it is assumed that \{x\_n\} converges, then lim x\_n =
\begin{enumerate}[label=(\alph*)]
    \item 3
    \item e
    \item sqrt(5)
    \item 0
\end{enumerate}
\item If c \textgreater{} 0 and f(x) = $e^x$ - cx for all real numbers x, then the minimum value of f is
\begin{enumerate}[label=(\alph*)]
    \item f(c)
    \item f($e^c$)
    \item f(1/c)
    \item f(log c)
\end{enumerate}
\item The set of integers Z with the binary operation "*" defined as a*b =a +b+ 1 for a, b in Z, is a group. The identity element of this group is
\begin{enumerate}[label=(\alph*)]
    \item 0
    \item 1
    \item -1
    \item 12
\end{enumerate}
\item Let y = f(x) be a solution of the differential equation x dy + (y - $xe^x$) dx = 0 such that y = 0 when x = 1. What is the value of f(2)?
\begin{enumerate}[label=(\alph*)]
    \item 1/(2 e)
    \item 1/e
    \item $e^2$/2
    \item 2 e
\end{enumerate}
\item Statement 1 | Any set of two vectors in $R^2$ is linearly independent. Statement 2 | If V = span(v 1, ... , vk) and \{v 1, ... , vk\} are linearly independent, then dim(V) = k.
\begin{enumerate}[label=(\alph*)]
    \item True, True
    \item False, False
    \item True, False
    \item False, True
\end{enumerate}
\end{enumerate}

\section*{True / False (10 marks)}
\textit{Answer True or False}
\begin{enumerate}[label=\arabic*.]
\setcounter{enumi}{5}
\item True or False: There are exactly three additive abelian groups of order 16 such that x + x + x + x = 0 for all x in the group.
\item True or False: The maximum possible order for an element of S\_6 is 6.
\item True or False: The nonzero integers under multiplication do not form a group because they lack an identity element.
\item True or False: The norm ||x||\_p can be expressed as an inner product specifically when p equals 2.
\item True or False: The largest order of an element in the group of permutations of 5 objects is 12.
\end{enumerate}

\section*{Long Form Response (20 marks)}
\textit{Answer each question with a clear and complete explanation.}
\begin{enumerate}[label=\arabic*.]
\setcounter{enumi}{10}
\item On her previous five attempts Sarah had achieved times, in seconds, of 86, 94, 97, 88 and 96, for swimming 50 meters. After her sixth try she brought her median time down to 92 seconds. What was her time, in seconds, for her sixth attempt?
\item Find the distance between the points (0,15) and (8,0).
\end{enumerate}

\vfill
\noindent\textit{End of Paper}
\end{document}